\subsection{分式环}

定理4. 设A是一个交换环，S是A的一个子集。设 \(\mathbf{A}_{\mathrm{S}}\) 为A的分式幺半群（仅配备乘法），其分母在S中（§2，第4号）。设 \(\varepsilon: \mathbf{A} \to \mathbf{A}_{\mathrm{S}}\) 为典范态射。在 \(\mathbf{A}_{\mathrm{S}}\) 上存在一种且仅一种加法，满足以下条件：

\begin{enumerate}
\def\labelenumi{(\alph{enumi})}
\item
  \(A_{s}\) ，连同这个加法和乘法，构成一个交换环；
\item
  \(\varepsilon\) 是一个环同态。
\end{enumerate}

假设已为 \(\mathbf{A}_{\mathrm{s}}\) 找到满足条件(a)和(b)的一个加法。

设 \(x, y \in A_s\) 。设 \(S'\) 是 \(A\) 的由 \(S\) 生成的稳定乘法子幺半群。存在 \(a, b \in A\) 和 \(p, q \in S'\) ，使得 \(x = a/p, y = b/q\) 。则

\[
x = \varepsilon (a q) \varepsilon (p q) ^ {- 1}, \quad y = \varepsilon (b p) \varepsilon (p q) ^ {- 1},
\]

从而

\begin{enumerate}
\def\labelenumi{(\arabic{enumi})}
\setcounter{enumi}{22}
\tightlist
\item
\end{enumerate}

\[
\begin{array}{r l} x + y & = (\varepsilon (a q) + \varepsilon (b p)) \varepsilon (p q) ^ {- 1} \\ & = \varepsilon (a q + b p) \varepsilon (p q) ^ {- 1} \\ & = (a q + b p) / p q. \end{array}
\]

这证明了加法的唯一性。

我们现在在 \(A_{s}\) 上通过设定 \(x + y = (aq + bp)/pq\) 来定义一个加法。有必要证明这个定义不依赖于a、b、p、q的选取。现在，若 \(a', b' \in A\) , \(p', q' \in S'\) 满足 \(x = a'/p', y = b'/q'\) ，则存在属于 \(S'\) 的 s 和 t，使得 \(ap's = a'ps\) , \(bq't = b'qt\) ，因此

\[
(a q + b p) \left(p ^ {\prime} q ^ {\prime}\right) (s t) = \left(a ^ {\prime} q ^ {\prime} + b ^ {\prime} p ^ {\prime}\right) (p q) (s t)
\]

从而

\[
(a q + b p) / p q = (a ^ {\prime} q ^ {\prime} + b ^ {\prime} p ^ {\prime}) / p ^ {\prime} q ^ {\prime}.
\]

容易验证 \(A_{s}\) 中的加法满足结合律和交换律，0/1是加法的单位元， \((-a)/p\) 是a/p的负元，并且 \(x(y + z) = xy + xz\) 对一切 \(x, y, z \in A_{s}\) 成立。若 \(a, b \in A\) ，则

\[
\varepsilon (a + b) = (a + b) / 1 = a / 1 + b / 1 = \varepsilon (a) + \varepsilon (b)
\]

从而 \(\varepsilon\) 是一个环同态。

定义8. 定理4中定义的环称为与S相伴的分式环，或分母在S中的分式环，并记为A{[}S \(^{-1}\) {]}.

\(\mathbf{A}[\mathbf{S}^{-1}]\) 的零元是 \(0 / 1\) ， \(\mathbf{A}[\mathbf{S}^{-1}]\) 的单位元是 \(1 / 1\) .

我们将在《交换代数》第二章，§ 2 中回到 \(\mathbf{A}[\mathrm{S}^{-1}]\) 的性质。

若S是A中可消去元素的集合，则环 \(A[S^{-1}]\) 称为A的全分式环。此时A被等同于 \(A[S^{-1}]\) 的一个子环，所借助的映射 \(\varepsilon\) 此时是单射的（I，§2，第4号，命题6）。

定理5. 设A是一个交换环，S是A的一个子集，B是一个环，f是A到B中的一个同态，使得 \(f(\mathbf{S})\) 的每个元素都可逆。存在一个且仅一个 \(\bar{f}\) ，它把 \(A[S^{-1}]\) 映入 B，使得 \(f = \bar{f} \circ \varepsilon\) .

我们知道（§2，第4号，定理1）存在一个且仅一个态射 \(\bar{f}\) ，它把乘法幺半群 A{[}S \(^{-1}\) {]} 映入乘法幺半群 B，使得 \(f = \bar{f} \circ \varepsilon\) 。设 \(a, b \in \mathbf{A}, p, q \in \mathbf{S}'\) （由S生成的A的稳定乘法子幺半群）。由于 \(f(A)\) 的元素两两交换，

\[
\begin{array}{r l} \bar {f} (a | \dot {p} + b | q) & = \bar {f} ((a q + b \dot {p}) | \dot {p} q) = f (a q + b \dot {p}) f (\dot {p} q) ^ {- 1} \\ & = (f (a) f (q) + f (b) f (\dot {p})) f (\dot {p}) ^ {- 1} f (q) ^ {- 1} \\ & = f (a) f (\dot {p}) ^ {- 1} + f (b) f (q) ^ {- 1} \\ & = \bar {f} (a | \dot {p}) + \bar {f} (b | q). \end{array}
\]

因此 \(\tilde{f}\) 是一个环同态。

